\documentclass[11pt]{article}

\usepackage[margin=1in]{geometry}
\usepackage{amsmath}
\usepackage{booktabs}
\usepackage{graphicx}
\usepackage{hyperref}
\usepackage{siunitx}
\usepackage{pgfplots}
\usepackage{tikz}
\usepgfplotslibrary{groupplots}
\pgfplotsset{compat=1.18}

\title{Reproducing LiteFNO on Gray-Scott Reaction--Diffusion}
\author{LiteFNO Reproduction Team}
\date{\today}

\begin{document}
\maketitle

\begin{abstract}
We report a focused reproduction of the Lightweight Fourier Neural Operator
(LiteFNO) on the Gray-Scott reaction--diffusion dataset from The Well benchmark.
Using the repository's default configuration for this task, the model reaches a
best validation RMSE of \num{0.013112} and validation VRMSE of \num{0.034148} at
epoch 29, with a test RMSE of \num{0.013201} and test VRMSE of \num{0.035189}.
The logged run contains 33 epochs (0--32) and a model size of 107,842 parameters.
We document the exact preprocessing, model, and training settings, provide
learning curves, and discuss limitations of the current run.
\end{abstract}

\section{Introduction and Motivation}
Neural operators provide a data-driven approach to learning solution operators
for time-dependent partial differential equations (PDEs). LiteFNO is a
lightweight variant designed to retain the modeling benefits of Fourier Neural
Operators while reducing parameter count and memory footprint. This paper
documents a targeted reproduction on the Gray-Scott reaction--diffusion system,
with emphasis on reproducibility and transparent reporting of configuration and
metrics.

Beyond benchmarking, efficient surrogate models matter for climate adaptation
workflows that must run under tight resource constraints. In sub-Saharan Africa,
smallholder farmers face climate variability that affects planting decisions,
water management, and crop resilience. Fast, low-footprint surrogate models can
help enable localized scenario analysis and rapid experimentation in applied
settings, where compute and connectivity are limited. This reproduction study
is a step toward validating lightweight operator learning as a building block
for such downstream decision-support systems.

\section{Dataset and preprocessing}
All datasets in this repository are drawn from The Well benchmark. The Gray-Scott
dataset is preprocessed to match the paper's downsampling and trajectory limits.
Table~\ref{tab:dataset} summarizes the exact preprocessing settings.

\begin{table}[t]
  \centering
  \caption{Gray-Scott preprocessing configuration.}
  \label{tab:dataset}
  \begin{tabular}{ll}
    \toprule
    Setting & Value \\
    \midrule
    Dataset name & \texttt{gray\_scott\_reaction\_diffusion} \\
    Downsample factor & 4 \\
    Max trajectories & 1000 \\
    Max steps & 60 \\
    Input steps & 1 \\
    Output steps & 1 \\
    Fields & 2 \\
    Stride & 1 \\
    Cache & \texttt{memory} \\
    Random seed & 0 \\
    Dataset key & \texttt{data} \\
    \bottomrule
  \end{tabular}
\end{table}

\section{Model and training setup}
LiteFNO uses a compact spectral architecture configured via YAML. The
Gray-Scott experiment inherits the base LiteFNO configuration and applies a
dataset-specific override for data paths and logging. Core settings are listed
in Table~\ref{tab:training}.

\begin{table}[t]
  \centering
  \caption{Model and training configuration for the Gray-Scott LiteFNO run.}
  \label{tab:training}
  \begin{tabular}{ll}
    \toprule
    Setting & Value \\
    \midrule
    Layers & 8 \\
    Width & 64 \\
    Rank & 32 \\
    Epochs (configured) & 200 \\
    Batch size & 512 \\
    Learning rate & 1e{-}3 \\
    LR step & 100 \\
    LR gamma & 0.5 \\
    Eval every & 1 \\
    Eval splits & train, valid \\
    Eval windows & [6, 12], [13, 30] \\
    Test at end & true \\
    Device & \texttt{cuda} \\
    Seed & 1337 \\
    AMP & true \\
    Checkpoint best metric & \texttt{valid\_vrmse} \\
    \bottomrule
  \end{tabular}
\end{table}

\section{Metrics}
We report root mean squared error (RMSE) and variance-normalized RMSE (VRMSE).
For prediction \(\hat{y}\) and target \(y\), the metrics are:
\begin{align}
\mathrm{RMSE} &= \sqrt{\mathbb{E}\left[(\hat{y} - y)^2\right]}, \\
\mathrm{VRMSE} &= \sqrt{\frac{\mathbb{E}\left[(\hat{y} - y)^2\right]}{\mathrm{Var}(y) + \epsilon}},
\end{align}
where \(\epsilon\) is a small stabilizer.

\section{Reproduction details}
The reproduction run uses the config
\texttt{configs/experiments/litefno\_gray\_scott\_reaction\_diffusion.yaml},
which inherits from \texttt{configs/experiments/base\_litefno.yaml}. Dataset
preprocessing settings are defined in
\texttt{configs/datasets/gray\_scott\_reaction\_diffusion.yaml}. The core
workflow is:
\begin{verbatim}
litefno download --config configs/datasets/gray_scott_reaction_diffusion.yaml
litefno preprocess --config configs/datasets/gray_scott_reaction_diffusion.yaml
litefno train --config configs/experiments/litefno_gray_scott_reaction_diffusion.yaml
litefno test --config configs/experiments/litefno_gray_scott_reaction_diffusion.yaml \
  --checkpoint outputs/checkpoints/gray_scott_reaction_diffusion/litefno/best.pt
\end{verbatim}
Metrics are logged as JSONL records to
\texttt{outputs/logs/gray\_scott\_reaction\_diffusion\_litefno.jsonl}. The
learning curves in Figure~\ref{fig:learning-curves} are generated from the
recorded log in this repository at
\texttt{logs/gray\_scott\_reaction\_diffusion\_litefno.jsonl}.

\section{Extensions \& Results}
This section reports the baseline reproduction results from the existing log
and outlines planned extensions for low-resource deployment.

The training log (\texttt{logs/gray\_scott\_reaction\_diffusion\_litefno.jsonl})
contains 33 epochs (0--32) plus a test evaluation recorded with step \(-1\).
Validation metrics improve rapidly early in training and plateau by the late
epochs. Table~\ref{tab:results} lists the best validation and test performance
and the logged parameter count.

\begin{table}[t]
  \centering
  \caption{Summary of Gray-Scott LiteFNO results.}
  \label{tab:results}
  \begin{tabular}{ll}
    \toprule
    Metric & Value \\
    \midrule
    Best validation epoch & 29 \\
    Best validation RMSE & \num{0.013112} \\
    Best validation VRMSE & \num{0.034148} \\
    Final logged epoch & 32 \\
    Final validation RMSE & \num{0.013310} \\
    Final validation VRMSE & \num{0.034442} \\
    Test RMSE & \num{0.013201} \\
    Test VRMSE & \num{0.035189} \\
    Parameters & 107,842 \\
    \bottomrule
  \end{tabular}
\end{table}

\begin{figure}[t]
  \centering
  \begin{tikzpicture}
    \begin{groupplot}[
      group style={group size=2 by 1, horizontal sep=1.2cm},
      width=0.48\textwidth,
      height=0.35\textwidth,
      xlabel={Epoch},
      grid=both,
      grid style={line width=.1pt, draw=gray!30}
    ]
      \nextgroupplot[ylabel={RMSE}]
      \addplot[blue] table [x=step, y=train_rmse, col sep=comma] {data/gray_scott_litefno_metrics.csv};
      \addplot[orange] table [x=step, y=valid_rmse, col sep=comma] {data/gray_scott_litefno_metrics.csv};
      \legend{Train,Valid}
      \nextgroupplot[ylabel={VRMSE}]
      \addplot[blue] table [x=step, y=train_vrmse, col sep=comma] {data/gray_scott_litefno_metrics.csv};
      \addplot[orange] table [x=step, y=valid_vrmse, col sep=comma] {data/gray_scott_litefno_metrics.csv};
      \legend{Train,Valid}
    \end{groupplot}
  \end{tikzpicture}
  \caption{Training and validation metrics for LiteFNO on Gray-Scott.}
  \label{fig:learning-curves}
\end{figure}

The gap between validation and test performance is small, suggesting good
generalization for this dataset under the configured preprocessing. The best
validation epoch (29) slightly outperforms the final epoch (32), indicating
modest metric oscillations late in training. Despite a compact parameter count,
LiteFNO achieves low VRMSE relative to the initial epoch, with validation RMSE
improving by roughly 70\% from epoch 0 to the best epoch.

Planned extensions in this repository include low-rank and width sweeps,
quantization, robustness checks, and explainability analyses. These extensions
are documented in the project roadmap but are not yet executed in the present
log, so no extension-specific results are reported here.

\section{Limitations}
This report is limited to a single LiteFNO run on one dataset and does not yet
include a baseline comparison (e.g., FNO-S) or multiple random seeds. The log
contains only 33 epochs, which is shorter than the 200 epochs configured in the
experiment, so the results should be interpreted as an early or partial run.
Windowed VRMSE metrics are not available in the current log, and the planned
extensions have not yet produced results.

\section{Conclusion}
We documented a LiteFNO reproduction on the Gray-Scott reaction--diffusion
dataset with full configuration transparency and quantitative results. The
current run delivers strong validation and test metrics with a compact model,
and provides a foundation for expanded sweeps and baseline comparisons.

\end{document}
